\documentclass[12pt]{article}
\usepackage{amsmath, amssymb, amsthm}
\usepackage{geometry}
\usepackage{hyperref}
 
\geometry{margin=1in}
 
\newtheorem*{theorem}{Theorem}
\theoremstyle{remark}
\newtheorem*{remark}{Remark}
\newtheorem*{note}{Note on the Proof Strategy}
 
\newcommand{\M}[1]{\mathcal{M}\!\left(#1\right)}
\newcommand{\Lc}{\mathcal{L}}
 
\title{\textbf{Axler Exercise 3C.12: Associativity of Matrix Multiplication}\\
\large A Proof and Conceptual Commentary}
\author{Linear Algebra Done Right, 4th Edition}
\date{}
 
\begin{document}
 
\maketitle
 
\section*{Context in the Text}
 
This is Exercise 12 in Section 3C (\emph{Matrices}) of \emph{Linear Algebra Done
Right}, 4th edition. Section 3C has just introduced matrix multiplication and
motivated its definition via composition of linear maps (Theorem 3.43). The
exercise asks the reader to prove associativity of matrix multiplication, and
explicitly invites a proof in the spirit of Emil Artin's remark:
 
\begin{quote}
\itshape
``It is my experience that proofs involving matrices can be shortened by 50\%
if one throws the matrices out.''
\end{quote}
 
\noindent The point of the exercise is to illustrate that matrix multiplication
was defined \emph{precisely so that} it mirrors composition of linear maps, and
that properties of matrices are best proved by leveraging that correspondence
rather than by index-chasing.
 
\section*{Prerequisite}
 
\begin{itemize}
    \item \textbf{Theorem 3.43 (Matrix of a product):} If $T \in \Lc(U, V)$ and
          $S \in \Lc(V, W)$, then
          \[
              \M{ST} = \M{S}\,\M{T}.
          \]
\end{itemize}
 
\section*{Statement}
 
\begin{theorem}[Exercise 3C.1 --- Associativity of matrix multiplication]
Suppose $A$, $B$, $C$ are matrices whose sizes are such that $(AB)C$ makes
sense. Then $A(BC)$ also makes sense and
\[
    (AB)C = A(BC).
\]
\end{theorem}
 
\section*{Why $A(BC)$ Makes Sense}
 
Suppose $C$ is an $m \times p$ matrix, $B$ is an $n \times m$ matrix, and $A$
is a $q \times n$ matrix, so that $AB$ is $q \times m$ and $(AB)C$ is
$q \times p$. Then $BC$ is $n \times p$ and $A(BC)$ is $q \times p$ as well.
Both products are defined and have the same size.
 
Equivalently, in operator language: if $C = \M{R}$, $B = \M{T}$, $A = \M{S}$
for linear maps $R \in \Lc(U, V)$, $T \in \Lc(V, W)$, $S \in \Lc(W, Z)$, then
both $(AB)C$ and $A(BC)$ represent the matrix of the composite
$S \circ T \circ R \in \Lc(U, Z)$, which is well defined.
 
\section*{Proof}
 
Every matrix is the matrix of some linear map with respect to standard bases.
So write $A = \M{S}$, $B = \M{T}$, $C = \M{R}$ for appropriate linear maps
$R$, $T$, $S$.
 
Operator composition is associative by definition --- $(ST)R$ and $S(TR)$ both
mean ``first apply $R$, then $T$, then $S$'' --- so:
\[
    (ST)R = S(TR).
\]
Applying $\mathcal{M}$ to both sides and using Theorem 3.43 twice on each side:
\[
    (AB)C = \M{ST}\,\M{R} = \M{(ST)R} = \M{S(TR)} = \M{S}\,\M{TR} = A(BC).
\]
\hfill$\blacksquare$
 
\section*{Commentary}
 
\begin{note}
The proof has two layers, and it is worth separating them clearly.
 
\begin{enumerate}
    \item \textbf{Operator composition is trivially associative.} There is
          nothing to prove here: $(ST)R$ and $S(TR)$ are the same function by
          definition. Associativity is a property inherited for free from the
          associativity of function composition.
 
    \item \textbf{The matrix correspondence $\mathcal{M}$ respects composition.}
          This is Theorem 3.43, which is non-trivial --- it is the entire
          motivation for why matrix multiplication is defined the way it is.
          Once you have it, any algebraic identity about composition of operators
          descends automatically to the same identity about their matrices.
\end{enumerate}
 
This is Artin's point in action. A direct proof by indices would require
writing out the $(i,k)$-entry of $(AB)C$ as
$\sum_{j}\bigl(\sum_{\ell} A_{i\ell}B_{\ell j}\bigr)C_{jk}$,
rearranging the sum using commutativity of scalar multiplication, and matching
it against the $(i,k)$-entry of $A(BC)$. That argument is correct but
unmotivating: it does not explain \emph{why} associativity holds, only that the
algebra works out. The operator proof, by contrast, makes the reason immediately
visible: matrices inherit associativity from functions, via $\mathcal{M}$.
\end{note}
 
\end{document}
